\subsection*{5.3. 条件 (A)}

\subsubsection*{5.3.1.}
设
\[
\rho:\widetilde X\longrightarrow\mathbb P^1
\]
为与某个线束相伴的“纤维化”, 该线束的轴与 $X$ 横截相交, 并假设其一般纤维光滑.

设 $U$ 为 $\mathbb P^1$ 的非空开子集, 使 $\rho$ 在其上光滑, 并设
\begin{equation}\tag{5.3.1.1}
\nu:U\hookrightarrow\mathbb P^1
\end{equation}
为包含映射. 由于层 $R^i\rho_*\mathbb Q_\ell(j)$ \underline{在 $U$ 上为扭常层}, 可见层
\[
\nu_*\nu^*R^i\rho_*\mathbb Q_\ell(j)
\]
本质上\underline{不依赖于} $U$ 的选取; 典范的“伴随”同态
\begin{equation}\tag{5.3.2}
R^i\rho_*\mathbb Q_\ell(j)
\longrightarrow
\nu_*\nu^*R^i\rho_*\mathbb Q_\ell(j)
\end{equation}
也同样不依赖于 $U$ 的选取.

\subsubsection*{5.3.3.}
注意, 对于\underline{任意}超平面截面线束, 弱 Lefschetz 定理蕴涵: 当 $i\leq n-2$ 时, 浸入
\[
\widetilde X\overset{(f\times\rho)}{\hookrightarrow}
X\times\mathbb P^1
\]
给出 $\mathbb P^1$ 上的层同构
\[
R^i\rho_*\mathbb Q_\ell(j)
\xleftarrow{\ \sim\ }
H^i(X,\mathbb Q_\ell(j))_{\mathbb P^1},
\qquad i\leq n-2.
\]
换言之, 当 $i\leq n-2$ 时, $R^i\rho_*\mathbb Q_\ell(j)$ 为常层, 因而当 $i\leq n-2$ 时, 映射 (5.3.2) 更是同构.

\subsubsection*{5.3.4.}
消灭圈的局部理论, 特别是\underline{Picard--Lefschetz 公式} [SGA 7, XV 3.4], 表明:
\underline{对 Lefschetz 线束, 当 $i\neq n$ 时同态 (5.3.2) 为同构, 而当 $i=n$ 时为满射}.

\subsubsection*{5.3.5.}
若\underline{所有}同态 (5.3.2) 都为同构 (等价地, 若取 $j=0$ 时对每个 $i$ 都成立), 则称 $\rho$ 或产生它的线束\underline{满足条件 (A)}.

我们将在下面 (6.3, 6.4) 看到, 条件 (A) 有很强的成立倾向.
